% !TEX encoding = UTF-8 Unicode
% !TeX program = xelatex
% !TEX TS-program = XeLaTeX

%\documentclass[12pt, a4paper]{article}
\documentclass[12pt, a5paper]{article}
\usepackage{mathtools}
\usepackage{unicode-math}
\usepackage{comment}
\usepackage{wrapfig}
\usepackage{xfrac}
\usepackage{tikz}
\usetikzlibrary{arrows, decorations.markings}

\usepackage[pass]{geometry}
%\geometry{legalpaper, landscape, margin=1in}
\usepackage{fullpage}
\usepackage{setspace}
\usepackage[localise]{xepersian}
\setstretch{1.195}
\settextfont[Scale=1]{XB Niloofar}
\setlatintextfont[Scale=1]{XB Niloofar}%{Junicode} %Junicode
\setdigitfont{XB Niloofar}%{XB Zar}

\setlength{\parindent}{0ex}

\definecolor{answercolor}{HTML}{0080ff}
\newcommand{\vsp}{\vspace{25pt}}
\newcommand{\پاسخ}[2]{#1}
%\thispagestyle{empty}
%\newcommand{\پاسخ}[2]{\color{answercolor}\textbf{پاسخ: }#2\color{black}\vspace{-0.3cm}}

\begin{document}
	\begin{tikzpicture}[remember picture,overlay]	\node[anchor=north,yshift=-1.5cm] at (current page.north) {\rl{به نام خدا}};
	\end{tikzpicture}%

\vspace{-1cm}
{\centering
	\textbf{کوئیز چهارم درس ساختمان داده‌ها و الگوریتم‌ها}\\
مدرس: مهدی صفرنژاد       - ۱۶ اردیبهشت ۹۸\\
\raggedleft\small
نام و نام خانوادگی: \\
}
\vspace{-0.5cm}
\پرش‌متوسط\پرخط‌ا\\
در تمام سؤالات پاسخ کوتاه (حدوداً بین یک تا دو خط) مدنظر است، مگر این که در سؤال به صراحت ذکر شده باشد.
\پاراگراف{سؤال ۱.}\sloppy
این جمله درست است یا غلط؟ با فرض درهم‌سازی یکنواخت ساده، احتمال این‌که در درهم‌سازی، سه عضو مشخص به یک خانۀ یکسان نگاشته شود برابر است با ${1}/{m^3}$ که $m$ اندازۀ جدول است. (در این سؤال به چگونگی حل مشکل برخورد کاری نداریم)

\پاسخ{\vsp}{
خیر، احتمال این که هر سه به یک خانۀ مشخص نگاشته شوند ${1}/{m^3}$ است و اگر خانۀ مشخصی مدنظر نباشد این احتمال برابر با ${1}/{m^2}$ است. 

}

\پاراگراف{سؤال ۲.}
با فرض درهم‌سازی یکنواخت ساده، یک جدول درهم‌سازی داریم که برای حل مشکل برخورد به جای ایجاد زنجیره (لیست پیوندی) از یک درخت دودویی جست‌وجوی متوازن استفاده می‌کنیم. زمان میانگین جست‌وجو در این جدول درهم‌سازی چقدر است؟ (ضریب بارگذاری $\alpha$ است. به حالتی که $\alpha=0$ است نیز توجه کنید.)

\پاسخ{\vsp}{

	
}

\پاراگراف{سؤال ۳.}
چرا وقتی از وارسی خطی برای رفع مشکل برخورد استفاده می‌کنیم، نمی‌توانیم یک عنصر را به راحتی از جدول حذف کنیم و خانۀ آن را خالی کنیم؟

\پاسخ{\vsp}{


}
	
\پاراگراف{سؤال ۴.}
جدول زیر یک جدول درهم‌سازی با وارسی خطی است که از تابع درهم‌سازی باقی‌مانده بر ۷ استفاده می‌کند. دو ترتیب مختلف از ورود اعضا را بنویسید که به این جدول منجر شود.
\begin{center}
\begin{tabular}{|c|c|c|c|c|c|c|}
	\hline
	0  & 1  & 2  & 3  & 4  & 5  & 6  \\\hline
	47 & 55 & 93 & 10 & 18 & 40 & 76\\\hline
\end{tabular}
\end{center}

\پاسخ{\vsp\vsp}{
	
}

\پاراگراف{سؤال ۵.}
در کلاس دیدیم اگر بخواهیم یک میلیون url را در یک مجموعه با فیلتر بلوم ذخیره کنیم به آرایه‌ای به طول نه میلیون و هفت تابع درهم‌سازی نیاز داریم. روش دیگر این است که از یک جدول درهم‌سازی به طول دو میلیون با وارسی خطی استفاده کنیم. برای هر روش یک برتری نسبت به روش دیگر ذکر کنید.


\پاسخ{\vsp\vsp}{
	
}

\پاراگراف{سؤال ۶.}
این جمله درست است یا غلط؟ در درهم‌سازی کوکو، زمان جست‌وجو \زیرخط{در بدترین حالت} $O(1)$ است.

\پاسخ{\vsp\vsp}{
درست است زیر 
	
}
\vspace{1pt}
\begin{flushleft}
موفق باشید
\end{flushleft}
\end{document}